%  Преамбула конспекта «Методы оптимизации»
%  Компилировать pdfLaTeX (кодировка utf8, шрифты T2A).

% Latin Modern не имеет T2A-начертаний, поэтому кириллица идёт через cm-super:
% семейство cmr в кодировке T2A.
\usepackage[T2A]{fontenc}
\usepackage[utf8]{inputenc}
\usepackage[english,russian]{babel}

\usepackage{amsmath,amssymb,amsthm,mathtools}
\usepackage{geometry}
\geometry{a4paper,margin=2.5cm}

\usepackage{tikz}
\usetikzlibrary{arrows.meta,patterns,decorations.pathmorphing,decorations.pathreplacing,shapes.geometric,calc,positioning}
\usepackage{pgfplots}
\usepgfplotslibrary{groupplots,fillbetween}
\pgfplotsset{compat=1.18}

\usepackage[most]{tcolorbox}
\tcbuselibrary{breakable}

\usepackage{enumitem}
\usepackage{booktabs}
\usepackage{hyperref}
\hypersetup{
  colorlinks=true,
  linkcolor=blue!55!black,
  citecolor=blue!55!black,
  urlcolor=blue!55!black,
  pdftitle={Методы оптимизации --- конспект}
}

%  Окружения для теорем
\theoremstyle{plain}
% Счётчики независимы: теоремы, леммы и утверждения нумеруются каждое своим рядом.
\newtheorem{theorem}{Теорема}
\newtheorem{proposition}{Утверждение}
\newtheorem{corollary}{Следствие}
\newtheorem*{corollary*}{Следствие}
\newtheorem{lemma}{Лемма}
\newtheorem{criterion}{Признак}
\newtheorem*{criterion*}{Признак}

\theoremstyle{definition}
\newtheorem{definition}{Определение}
\newtheorem{example}{Пример}
\newtheorem{exercise}{Упражнение}

\theoremstyle{remark}
\newtheorem*{remark}{Замечание}
\newtheorem*{reminder}{Напоминание}

\renewcommand{\qedsymbol}{$\blacksquare$}

%  Пометки конспекта
%    \begin{idea} ... \end{idea} — пояснения, которых не было на лекции
%    \begin{warn} ... \end{warn} — неразборчиво / описка / требует сверки
\newtcolorbox{idea}[1][]{
  breakable,
  enhanced,
  colback=blue!4!white,
  colframe=blue!45!black,
  boxrule=0.4pt,
  left=6pt,right=6pt,top=5pt,bottom=5pt,
  title={Пояснение},
  fonttitle=\bfseries\small,
  coltitle=white,
  attach boxed title to top left={xshift=6pt,yshift=-2pt},
  boxed title style={colback=blue!45!black,boxrule=0pt,arc=1pt},
  #1
}

\newtcolorbox{warn}[1][]{
  breakable,
  enhanced,
  colback=orange!6!white,
  colframe=orange!65!black,
  boxrule=0.4pt,
  left=6pt,right=6pt,top=5pt,bottom=5pt,
  title={Сверить с записью},
  fonttitle=\bfseries\small,
  coltitle=white,
  attach boxed title to top left={xshift=6pt,yshift=-2pt},
  boxed title style={colback=orange!65!black,boxrule=0pt,arc=1pt},
  #1
}

%  Пошаговая запись метода: \begin{algo}[title={Метод ...}] ... \end{algo}
\newtcolorbox{algo}[1][]{
  breakable,
  enhanced,
  colback=black!3!white,
  colframe=black!60,
  boxrule=0.4pt,
  left=6pt,right=6pt,top=5pt,bottom=5pt,
  title={Алгоритм},
  fonttitle=\bfseries\small,
  coltitle=white,
  attach boxed title to top left={xshift=6pt,yshift=-2pt},
  boxed title style={colback=black!60,boxrule=0pt,arc=1pt},
  #1
}

% Шаги внутри algo
\newlist{steps}{enumerate}{1}
\setlist[steps]{label=\textbf{Шаг \arabic*.}, leftmargin=*, labelsep=0.6em,
                topsep=4pt, itemsep=3pt}

% Отметка начала лекции
\newcommand{\lecture}[2]{%
  \bigskip
  \noindent\hfill\textit{\small #1\ifx\relax#2\relax\else, #2\fi}\\[-0.3em]
  \noindent\rule{\textwidth}{0.4pt}
  \medskip
}

%  Макросы
\newcommand{\R}{\mathbb{R}}
\newcommand{\NN}{\mathbb{N}}
\newcommand{\eps}{\varepsilon}
\newcommand{\T}{^{\mathsf{T}}}
\newcommand{\rank}{\operatorname{rank}}
\newcommand{\sgn}{\operatorname{sign}}
\newcommand{\norm}[1]{\left\lVert #1 \right\rVert}
\newcommand{\abs}[1]{\left\lvert #1 \right\rvert}
\DeclareMathOperator*{\argmin}{arg\,min}
\DeclareMathOperator*{\argmax}{arg\,max}

% Индексные множества канонической задачи.
%   N — индексы переменных, M — индексы ограничений,
%   N_0^+ / N_0^0 — положительные / нулевые компоненты допустимой точки,
%   N^+ — базисное множество (пополненное до |M|).
\newcommand{\Np}{N^{+}}
\newcommand{\Nz}{N^{0}}
\newcommand{\Nzp}{N_0^{+}}
\newcommand{\Nzz}{N_0^{0}}

% Сокращение «опорный вектор», как в тетради.
\newcommand{\ov}{о/в}

% Клетка транспортной таблицы: #1 — тариф c_ij (мелко, сверху),
% #2 — объём перевозки x_ij (обычным кеглем, снизу).
\newcommand{\tc}[2]{\shortstack[r]{\scriptsize $#1$\\ $#2$}}
% Пустая (незаполненная) клетка таблицы.
\newcommand{\tcx}[1]{\tc{#1}{\times}}

% Маркер «++ в кружке» из тетради: доказательство будет дописано позже.
\newcommand{\plusplus}{%
  \tikz[baseline=(pp.base)]{%
    \node[draw, ellipse, inner xsep=3pt, inner ysep=1pt, line width=0.4pt] (pp)
      {\footnotesize $+\,+$};}%
}

%  Быстрая сборка: \showfiguresfalse выключает все TikZ-рисунки.
\newif\ifshowfigures
\showfigurestrue

\newcommand{\inputfig}[1]{%
  \ifshowfigures
    \input{#1}%
  \else
    \begin{center}\fbox{\small рисунок выключен: \texttt{\detokenize{#1}}}\end{center}%
  \fi
}
